\documentclass[10pt,a4paper]{amsart}
\usepackage[margin=19mm]{geometry}
\usepackage{amsmath,amssymb,mathtools}
\usepackage[scr=rsfs,cal=euler]{mathalfa}
\usepackage{xfrac}
\usepackage[unicode,hidelinks,bookmarks=false]{hyperref}
\newcommand{\ie}{i.e.\ }
\newcommand{\cf}{cf.\ }
\newcommand{\resp}{respectively\ }
\newcommand{\Subsection}{\subsection}
\providecommand{\nobreakdash}{\nobreak\hbox{-}}
\setlength{\emergencystretch}{2em}
\multlinegap0pt
\usepackage{amssymb}
\newtheorem{theorem}{Theorem}
\newtheorem{proposition}{Proposition}
\DeclareMathOperator{\GL}{\mathrm{GL}}
\DeclareMathOperator{\PGL}{\mathrm{PGL}}
\DeclareMathOperator{\PGU}{\mathrm{PGU}}
\DeclareMathOperator{\PSL}{\mathrm{PSL}}
\DeclareMathOperator{\PSU}{\mathrm{PSU}}
\title{Simple Lie Groups of type $A_n$ as Galois groups over $\mathbf{Q}$}
\author{Stepan Nesterov}
\date{Source: arXiv:2604.27402v1 (CC BY 4.0)}
\begin{document}
\maketitle

\noindent\emph{Source and licence.} Simple Lie Groups of type $A_n$ as Galois groups over $\mathbf{Q}$. Version arXiv:2604.27402v1 (CC BY 4.0), distributed by arXiv under the Creative Commons Attribution 4.0 International licence, http://creativecommons.org/licenses/by/4.0/, as recorded in the arXiv licence metadata for this exact version and shown on the abstract page https://arxiv.org/abs/2604.27402v1. Full licence notice: \texttt{licenses/arxiv-2604.27402-CC-BY-4.0.txt}.\par\smallskip

\noindent\textit{Editorial scope.} Counted unit: the article's theorem final (source0.tex lines 52--58). The article's complete original proof of it is delivered with it (source0.tex lines 219--223, one \texttt{proof} environment of 5 lines, which the article places away from the statement). No mathematics is written, repaired or re-derived: the statement and the proof are the article's own lines, in the article's own notation, and no retained line is substituted or re-flowed.  Retained uncounted as the source-local context the proof needs: lines 99--105; lines 146--148; lines 158--160; lines 166--168.  Nothing was omitted from any retained span, so the package carries no omission record.  No retained line contains a citation, so the wrapper carries no bibliography and no bibliography entry.  No retained line contains a figure environment, an image-inclusion command or drawing code, so no image is delivered and the package declares no asset dependency.  Editorial formatting only, in addition: an AMS \texttt{amsart} preamble that restates the article's own declarations and macros used by the retained lines, namely the environments \texttt{theorem}, \texttt{proposition} on separate counters, together with the standard packages those lines need.  The retained environments print in this document as Theorem 1, Proposition 1 in order of appearance, whereas the article prints the same items with the numbering of its own full document.  The complete native source of the article as distributed in the arXiv e-print for this exact version is bundled unchanged as \texttt{sources/199469/source0.tex}.
	\begin{theorem} \label{main}
		Let $l$ be an odd prime, let $\mu$ be an integer dividing $l-1$, let $n$ be an integer, such that $n \equiv -2 \mod \mu$, and let $p$ be an odd prime, such that $p \neq l$.  
		\begin{itemize}
			\item If $\frac{l-1}{\mathrm{ord (p \mod l)}}$ is odd, then one can realize a group between  $\PGU(n,q)$ or $\PSU(n,q)$ with $q=p^{\frac{\mathrm{ord} (p \mod l)}{2}}$, as a Galois group over the field $E_{\mu,l}$.
			\item If $\frac{l-1}{\mathrm{ord (p \mod l)}}$ is even, then one can realize a group between  $\PGL(n,q)$ or $\PSL(n,q)$ with $q=p^{\mathrm{ord} (p \mod l)}$, as a Galois group over the field $E_{\mu,l}$.
		\end{itemize}
	\end{theorem}

 	\begin{theorem} \label{unitary}
 		For the projective Galois representations $G_{\mathbf{Q}} \to \PGU(n,q)$ constructed in the previous theorem, if  $l {\not |}\ n$ and if either $q$ is square or $\gcd(n,q+1)=\gcd(\frac{n}{2},q+1)$, then the image is either $\PSU(n,q)$ or $\PSU(n,q)$ has index two in the image.
 	\end{theorem}

	\begin{theorem} \label{linear}
		For the projective Galois representations $G_{\mathbf{Q}} \to \PGL(n,q)$ constructed in the previous theorem, if  $l {\not |}\ n$, then the image is either $\PSL(n,q)$ or $\PSL(n,q)$ has index two in the image.
	\end{theorem}

	\begin{proposition} \label{odd'}
		For the Galois representations $G_{\mathbf{Q}} \to \GL(n,q)$ constructed in theorem \ref{main}, when we write $\det \rho = \varepsilon^{\frac{n}{2}} \psi$, the character $\psi$ has odd order.
	\end{proposition}

	\begin{theorem} \label{final}
		Let $l$ be an odd prime, and let $n$ be a natural number, not congruent to $-2$ modulo $l$. Let $e = \mathrm{ord} \ p \mod l$. Then:
		\begin{itemize}
			\item If $\frac{l-1}{e}$ is even, then $\PSL((l-1)n-2,p^e)$ is a Galois group over $\mathbf{Q}$.
			\item If $\frac{l-1}{e}$ is odd and, in addition to all the previous assumptions, either $e$ is a multiple of $4$ or if $2^r$ is the largest power of $2$ dividing $p+1$, then $\frac{l-1}{2}n \equiv 1 \mod 2^r$, then $\PSU((l-1)n - 2, p^{\frac{e}{2}})$ is a Galois group over $\mathbf{Q}$.
		\end{itemize}
	\end{theorem}

	\begin{proof} (of theorem \ref{final})
		Let $l$ be an odd prime, let $n$ be a natural number, not congruent to $-1$ mod $l$. Let $e = \mathrm{ord}\ p \mod l$. Consider the Galois representations constructed in theorem \ref{main} with $\mu = l-1$ and $n = (l-1)n - 2$. If $\frac{l-1}{e}$ is even, it follows from the propositions \ref{linear} and \ref{odd'} that these Galois representations have projective image $\PSL(n,p^e)$. \par 
		If $\frac{l-1}{e}$ is odd, we sometimes have to verify extra congruence conditions. If $e$ is a multiple of $4$, then $p^{\frac{e}{2}}$ is a square, and \ref{unitary} applies. Otherwise, let $r = v_2(p+1)$, and assume that $\frac{l-1}{2} n \equiv 1 \mod 2^r$. I claim that for $q=p^{\frac{e}{2}}$, we have $\gcd((l-1)n-2,q+1)=\gcd(\frac{(l-1)n-2}{2},q+1)$, so that \ref{unitary} applies. Indeed, the only case when $\gcd((l-1)n-2,q+1)>\gcd(\frac{(l-1)n-2}{2},q+1)$ is when $v_2(q+1) > v_2(\frac{(l-1)n-2}{2})= v_2(\frac{l-1}{2}n - 1)$. I claim that $v_2(q+1) = v_2(p+1)$, so this inequality does not happen. \par 
		To prove this claim, recall that we are in the case where $\frac{e}{2}$ is odd. By assumtion, $p + 1 \equiv 2^r \mod 2^{r+1}$. From this, we obtain $p^{\frac{e}{2}} \equiv (2^r - 1)^{\frac{e}{2}} \equiv (-1)^{\frac{e}{2}} + \frac{e}{2} 2^r \equiv -1 + 2^r \mod 2^{r+1}$, so $p^{\frac{e}{2}} + 1$ is also not divisible by $2^{r+1}$.
	\end{proof}

\end{document}
